% ============================================================================
%  面向 SCADA 工控系统的轻量化入侵检测方法研究
%  硕士论文模板 - 教学诊断版（含 \pitfall{} 与 \todo{} 标注）
%
%  本模板面向指导老师：在学生易错位置预埋诊断标记。
%  - \pitfall{...}    红色警示框，描述"学生常犯错误"，老师可直接对照评阅
%  - \todo[owner]{...} 黄色待办框，标记需学生补充的内容
%  - \grade[level]{...} 老师评分用，level=A/B/C/D
%
%  编译: xelatex main -> biber main -> xelatex main -> xelatex main
%  编译开关: \finaltrue  = 最终版(隐藏所有警示);
%           \finalfalse = 草稿版(显示警示,默认)
% ============================================================================

\documentclass[12pt,a4paper,oneside,openany]{ctexart}

\newif\iffinal
\finalfalse   % 草稿模式，老师可改为 \finaltrue 看最终版样式

% ---------- 数学与符号 ----------
\usepackage{amsmath,amssymb,amsthm,amsfonts,bm}
\usepackage{mathrsfs}

% ---------- 图形与表格 ----------
\usepackage{graphicx}
\usepackage{booktabs}
\usepackage{array}
\usepackage{multirow}
\usepackage{multicol}
\usepackage{makecell}
\usepackage{tabularx}
\usepackage{longtable}
\usepackage{adjustbox}
\usepackage{caption}
\usepackage{subcaption}
\usepackage{colortbl}
\usepackage{xcolor}
\usepackage{tcolorbox}
\tcbuselibrary{skins,breakable}

% ---------- 代码与算法 ----------
\usepackage{listings}
\usepackage{algorithm}
\usepackage{algorithmic}
\usepackage{float}

% ---------- 引用与链接 ----------
\usepackage[colorlinks=true,
            linkcolor=black,
            citecolor=blue,
            urlcolor=blue,
            bookmarksdepth=4]{hyperref}

% ---------- 参考文献（BibLaTeX + GB/T 7714） ----------
\usepackage[backend=biber, style=gb7714-2015, gbpub=false, gbnamefmt=lowercase]{biblatex}
\addbibresource{references/refs.bib}

% ---------- 页面布局 ----------
\usepackage{geometry}
\geometry{a4paper,left=2.5cm,right=2.5cm,top=2.5cm,bottom=2.5cm,
          headheight=15pt,headsep=0.4cm,footskip=0.8cm}

% ---------- 页眉页脚 ----------
\usepackage{fancyhdr}
\pagestyle{fancy}
\fancyhf{}
\fancyhead[L]{\small\leftmark}
\fancyhead[R]{\small\rightmark}
\fancyfoot[C]{\thepage}
\renewcommand{\headrulewidth}{0.4pt}

% ---------- 章节标题格式 ----------
\ctexset{
  chapter={
    name={第,章}, number=\chinese{chapter},
    format={\centering\heiti\zihao{2}},
    nameformat={\heiti\zihao{2}}, titleformat={\heiti\zihao{2}},
    beforeskip={1.0ex}, afterskip={1.5ex}, pagestyle=fancy
  },
  section={
    format={\heiti\zihao{3}}, number=\chinese{section},
    beforeskip={0.8ex}, afterskip={0.4ex}
  },
  subsection={
    format={\heiti\zihao{4}}, number=\arabic{subsection},
    beforeskip={0.6ex}, afterskip={0.3ex}
  }
}

% ---------- 定理环境 ----------
\newtheoremstyle{mystyle}{1.5ex}{1.5ex}{\kaishu}{}{\heiti}{：}{0.5em}{}
\theoremstyle{mystyle}
\newtheorem{definition}{定义}[chapter]
\newtheorem{theorem}{定理}[chapter]
\newtheorem{lemma}{引理}[chapter]
\newtheorem{proposition}{命题}[chapter]
\newtheorem{corollary}{推论}[chapter]
\newtheorem{example}{例}[chapter]
\renewcommand{\proofname}{证明}

% ---------- 列表 ----------
\usepackage{enumitem}
\setlist{nosep, leftmargin=*}

% ---------- 代码样式 ----------
\definecolor{codebg}{rgb}{0.97,0.97,0.97}
\definecolor{codekw}{rgb}{0.0,0.0,0.8}
\definecolor{codecomment}{rgb}{0.25,0.5,0.35}
\definecolor{codestring}{rgb}{0.6,0.1,0.1}
\lstset{
  basicstyle=\ttfamily\small, backgroundcolor=\color{codebg},
  keywordstyle=\color{codekw}\bfseries,
  commentstyle=\color{codecomment}\itshape,
  stringstyle=\color{codestring}, numbers=left,
  numberstyle=\tiny\color{gray}, numbersep=8pt,
  frame=single, framerule=0.4pt, rulecolor=\color{gray!50},
  breaklines=true, showstringspaces=false, tabsize=2, captionpos=b
}

% ============================================================================
%                    *** 教学诊断宏（核心） ***
% ============================================================================
% 这是模板的灵魂：老师通过这些宏命令精准定位学生问题
%
% \pitfall{标题}{描述}
%   红色警示框，描述"学生常犯错误"
%
% \todo[负责人]{描述}
%   黄色待办框，标记需要补充的内容
%
% \grade[等级]{评语}
%   老师评分块（草稿版不显示）
% ============================================================================

% --- Pitfall 宏（学生常犯错误） ---
\newcommand{\pitfall}[2]{%
  \iffinal
    % 最终版隐藏
  \else
    \begin{tcolorbox}[
      enhanced, breakable,
      colback=red!5!white, colframe=red!75!black,
      title={\textbf{⚠️ 学生常见错误 [#1]}},
      fonttitle=\bfseries\heiti\zihao{-5},
      boxrule=0.6pt, arc=2pt, left=6pt, right=6pt,
      top=4pt, bottom=4pt,
      fontupper=\kaishu\zihao{-5}
    ]
    #2
    \end{tcolorbox}
  \fi
}

% --- Todo 宏（待补充内容） ---
\newcommand{\todo}[2][学生]{%
  \iffinal
  \else
    \begin{tcolorbox}[
      enhanced, breakable,
      colback=yellow!10!white, colframe=orange!80!black,
      title={\textbf{📝 待补充（#1）}},
      fonttitle=\bfseries\heiti\zihao{-5},
      boxrule=0.6pt, arc=2pt, left=6pt, right=6pt,
      top=4pt, bottom=4pt,
      fontupper=\kaishu\zihao{-5}
    ]
    #2
    \end{tcolorbox}
  \fi
}

% --- Grade 宏（老师评分批注） ---
\newcommand{\grade}[2][B]{%
  \iffinal
  \else
    \begin{tcolorbox}[
      enhanced, breakable,
      colback=blue!5!white, colframe=blue!60!black,
      title={\textbf{✏️ 老师评语 [等级：#1]}},
      fonttitle=\bfseries\heiti\zihao{-5},
      boxrule=0.6pt, arc=2pt, left=6pt, right=6pt,
      top=4pt, bottom=4pt,
      fontupper=\kaishu\zihao{-5}
    ]
    #2
    \end{tcolorbox}
  \fi
}

% --- 优秀示例标记 ---
\newcommand{\good}[1]{%
  \iffinal
  \else
    \textcolor{green!50!black}{\textbf{[✓ 范例]}} #1
  \fi
}

% ---------- 常用宏 ----------
\newcommand{\figref}[1]{图~\ref{#1}}
\newcommand{\tabref}[1]{表~\ref{#1}}
\newcommand{\eqref}[1]{式~(\ref{#1})}
\newcommand{\chapref}[1]{第~\ref{#1}~章}
\newcommand{\secref}[1]{第~\ref{#1}~节}
\newcommand{\Fone}{F_1}
\newcommand{\MacroFone}{F_1^{\text{macro}}}
\newcommand{\WeightedFone}{F_1^{\text{weighted}}}
\newcommand{\PRAUC}{\text{PR-AUC}}
\newcommand{\ROCAUC}{\text{ROC-AUC}}

% ---------- 标题元信息 ----------
\title{
  \vspace{-1cm}
  \heiti\zihao{2}面向 SCADA 工控系统的\\
  \vspace{0.3cm}
  轻量化入侵检测方法研究\\[0.5cm]
  \fangsong\zihao{4}——基于时序卷积网络与极限压缩量化
}
\author{}
\date{}

% ============================================================================
\begin{document}

% ---------- 封面 ----------
% 封面页
\begin{titlepage}
\begin{center}

\vspace*{2cm}

{\heiti\zihao{1} 硕士学位论文}\\[0.3cm]
{\fangsong\zihao{4} （学术学位）}\\[2.5cm]

\begin{minipage}{0.85\textwidth}
\centering
{\heiti\zihao{2} 面向 SCADA 工控系统的}\\[0.4cm]
{\heiti\zihao{2} 轻量化入侵检测方法研究}\\[0.5cm]
{\fangsong\zihao{4} ——基于时序卷积网络与极限压缩量化}
\end{minipage}

\vspace{3cm}

\begin{tabular}{rl}
  {\fangsong\zihao{4} 研\hspace{2em}究\hspace{2em}生：}  & {\fangsong\zihao{4} \underline{\hspace{5cm}}} \\[0.6cm]
  {\fangsong\zihao{4} 导\hspace{2em}师：}        & {\fangsong\zihao{4} \underline{\hspace{5cm}}} \\[0.6cm]
  {\fangsong\zihao{4} 学科专业：}              & {\fangsong\zihao{4} \underline{\hspace{5cm}}} \\[0.6cm]
  {\fangsong\zihao{4} 研究方向：}              & {\fangsong\zihao{4} \underline{\hspace{5cm}}} \\[0.6cm]
\end{tabular}

\vfill

{\fangsong\zihao{4} 二〇二六年六月}

\end{center}
\end{titlepage}

% 中文扉页
\thispagestyle{empty}
\begin{center}
\vspace*{2cm}

{\fangsong\zihao{4} 分类号：\underline{\hspace{2cm}}\hspace{1cm} 密级：\underline{\hspace{2cm}}}\\[0.5cm]
{\fangsong\zihao{4} UDC：\underline{\hspace{2cm}}\hspace{1cm} 编号：\underline{\hspace{4cm}}}\\[4cm]

\begin{minipage}{0.9\textwidth}
\centering
{\heiti\zihao{2} 面向 SCADA 工控系统的}\\[0.4cm]
{\heiti\zihao{2} 轻量化入侵检测方法研究}
\end{minipage}

\vspace{3cm}

{\fangsong\zihao{4} 学位申请人：\underline{\hspace{5cm}}}\\[0.6cm]
{\fangsong\zihao{4} 指导教师：\underline{\hspace{5cm}}}\\[0.6cm]
{\fangsong\zihao{4} 学科专业：\underline{\hspace{5cm}}}\\[0.6cm]

\vfill

{\fangsong\zihao{4} 二〇二六年六月}
\end{center}

\clearpage


% ---------- 诚信声明 ----------
% 诚信声明
\thispagestyle{empty}
\begin{center}
{\heiti\zihao{2} 学位论文原创性声明}
\end{center}

\vspace{0.5cm}

{\fangsong\zihao{4} 本人郑重声明：所呈交的学位论文，是本人在导师的指导下，独立
进行研究工作所取得的成果。除文中已经注明引用的内容外，本论文不含任何其他
个人或集体已经发表或撰写过的作品成果。对本文的研究做出重要贡献的个人和集
体，均已在文中以明确方式标明。本人完全意识到本声明的法律结果由本人承担。}

\vspace{0.5cm}

\begin{flushright}
{\fangsong\zihao{4} 学位论文作者签名：\underline{\hspace{3cm}}}\\[0.3cm]
{\fangsong\zihao{4} 签字日期：\underline{\hspace{6cm}}}
\end{flushright}

\vspace{1.5cm}

\begin{center}
{\heiti\zihao{2} 学位论文版权使用授权书}
\end{center}

\vspace{0.5cm}

{\fangsong\zihao{4} 本学位论文作者完全了解学校有关保留、使用学位论文的规定，
同意学校保留并向国家有关部门或机构送交论文的复印件和电子版，允许论文被查阅
和借阅。本人授权学校可以将学位论文的全部或部分内容编入有关数据库进行检索，
可以采用影印、缩印或扫描等复制手段保存和汇编本学位论文。}

\vspace{0.5cm}

{\fangsong\zihao{4} 保密的学位论文在解密后适用本授权书。}

\vspace{1.5cm}

\begin{tabular}{rl}
{\fangsong\zihao{4} 学位论文作者签名：} & {\fangsong\zihao{4} \underline{\hspace{3cm}}} \\[0.6cm]
{\fangsong\zihao{4} 导师签名：}         & {\fangsong\zihao{4} \underline{\hspace{3cm}}} \\[0.6cm]
{\fangsong\zihao{4} 签字日期：}         & {\fangsong\zihao{4} \underline{\hspace{5cm}}} \\[0.6cm]
\end{tabular}

\clearpage


% ---------- 中文摘要 ----------
% 中文摘要
\chapter*{\hfill 中 文 摘 要 \hfill}
\addcontentsline{toc}{chapter}{中文摘要}
\thispagestyle{fancy}

\setcounter{page}{1}
\pagenumbering{Roman}

\par SCADA（监控与数据采集）系统作为国家关键基础设施的神经中枢，广泛应用于燃气、
电力、水务等工业控制领域，其安全性直接关系到公共安全与经济稳定。随着工业互联网
与人工智能技术的深度融合，传统基于规则与签名的入侵检测方法已难以应对日益复杂的
网络攻击，亟需研究面向 SCADA 系统的智能化、轻量化入侵检测方法。

\par 然而，SCADA 工控场景对入侵检测模型提出了独特而严苛的要求：一方面，工业控
制协议（如 Modbus）的时序结构蕴含丰富的领域语义，通用时序模型难以充分利用；
另一方面，受限于边缘设备的算力与存储，检测模型必须兼顾精度与部署可行性。现有
研究多沿用通用网络入侵检测的特征体系，缺乏协议层语义编码，且鲜有工作系统性地
探索从算法设计到芯片部署的完整优化链条。

\par 针对上述问题，本文以 IanArffDataset 数据集（274{,}000 条 SCADA 燃气管道记
录，20 维原始特征，三层标签）为研究对象，围绕``轻量化入侵检测''这一核心目标，
开展了从数据特征、模型架构到边缘部署的全栈研究，主要贡献与创新点如下：

\par （1）提出了一种基于 SCADA 协议语义的领域特征工程方法。深入分析 Modbus 协
议的 4 行主-从交互结构，将 4 行视为一个 Modbus 周期，在 17 维最小特征的基础上
设计 27 维协议语义特征（3 行级 + 8 窗口级）。在 LightGBM、TCN 等多个模型上的对
比实验表明，所提特征集在所有模型上一致优于原始 44 维特征，验证了``领域特征工
程的边际收益高于模型升级''的结论。

\par （2）提出了一种融合 SE 注意力的 TCN 轻量入侵检测模型（TCN-SE）。在系统
对比 27 种典型时序模型的基础上，发现 TCN 因果卷积的归纳偏置天然适配 Modbus
短时序结构；进一步引入 SE 通道注意力，仅增加 3{,}300 参数（+4.3\%）即在测试
集上取得 Macro-F1 = 0.8340，PR-AUC = 0.9025，首次在该任务上实现深度学习模型对
梯度提升树（LightGBM，Macro-F1 = 0.8337）的超越。通过窗口长度（$w=4, 8, 16$）
与通道宽度（ch=8 $\sim$ 64）的消融实验，揭示了``8 步窗口（2 个 Modbus 周期）
是该数据集的 Pareto 最优结构''，并验证了 SE 模块对少数类召回率的提升作用。

\par （3）提出了一种面向 MCU 部署的极限压缩与混合精度量化方案。通过系统的通道
宽度消融，得到参数量仅 4{,}237 的 TCN V19 变体，PR-AUC 反升至 0.9133（+0.67\%），
挑战了``模型越大性能越好''的传统认知。进一步提出``权重 INT8 + 激活 FP32''的
混合精度量化方案，在 Cortex-M4 168\,MHz 平台上实现 5.5\,KB 模型、0.47\,ms 推理
延迟，量化前后预测一致率达 99.80\%，SNR $\geq$ 38\,dB 无损。完整 380 行 C99 推理
代码已开源，可直接部署于 STM32、ESP32、Arduino 等嵌入式平台。

\par 本文围绕``数据-算法-芯片''三层递进的研究路线，形成了从领域特征到 MCU 部署
的完整创新链条。三个创新点在公开数据集上取得了 Macro-F1 = 0.8340（单模）与
PR-AUC = 0.9133（压缩单模）的最优结果，并验证了在 5.5\,KB 嵌入式模型上保持工
业级检测性能的可行性，为 SCADA 工控入侵检测的工程化落地提供了可复现的参考方
案。

\vspace{1em}

\noindent\textbf{关键词：} SCADA；入侵检测；时序卷积网络；模型压缩；混合精度量化；嵌入式部署

\clearpage


% ---------- 英文摘要 ----------
% 英文摘要
\chapter*{\hfill ABSTRACT \hfill}
\addcontentsline{toc}{chapter}{ABSTRACT}

\par SCADA (Supervisory Control and Data Acquisition) systems serve as the central
nervous system of critical national infrastructure, widely deployed in gas, power,
and water industrial control scenarios. Security breaches directly threaten public
safety and economic stability. With the deep integration of Industrial Internet of
Things (IIoT) and artificial intelligence, traditional signature-based intrusion
detection methods can no longer cope with increasingly sophisticated cyber attacks.
There is an urgent need to develop intelligent, lightweight intrusion detection
methods tailored for SCADA systems.

\par However, SCADA scenarios impose unique and stringent requirements on detection
models. On one hand, the temporal structure of industrial control protocols (e.g.,
Modbus) contains rich domain semantics that general time-series models fail to
exploit. On the other hand, constrained by the limited computing power and storage
of edge devices, detection models must balance accuracy with deployability. Existing
work often re-uses generic network intrusion detection features without protocol-level
semantic encoding, and few studies systematically explore the complete optimization
chain from algorithm design to on-chip deployment.

\par To address these issues, this thesis takes the IanArffDataset (274{,}000
SCADA gas-pipeline records, 20 raw features, three-layer labels) as the research
object, and conducts a full-stack study covering data features, model architecture,
and edge deployment, all centered on the core goal of ``lightweight intrusion
detection''. The main contributions are as follows:

\par (1) A domain feature engineering method based on SCADA protocol semantics is
proposed. By analyzing the 4-row master-slave interaction structure of Modbus
and treating 4 rows as one Modbus cycle, 27-dimensional protocol semantic features
(3 row-level + 8 window-level) are designed on top of the 17-dimensional minimal
feature set. Comparison experiments on multiple models (LightGBM, TCN, etc.)
demonstrate that the proposed feature set consistently outperforms the original
44-dimensional features, validating the conclusion that ``the marginal benefit of
domain feature engineering exceeds that of model upgrading''.

\par (2) A lightweight intrusion detection model (TCN-SE) that integrates SE
attention into Temporal Convolutional Network is proposed. Through a systematic
comparison of 27 typical time-series models, the inductive bias of TCN's causal
convolution is found to naturally fit Modbus short-term temporal structure.
Further introducing the SE channel attention, with only 3{,}300 additional parameters
(+4.3\%), the model achieves Macro-F1 = 0.8340 and PR-AUC = 0.9025 on the test
set, marking the first time a deep learning model surpasses gradient boosting
trees (LightGBM, Macro-F1 = 0.8337) on this task. Ablation studies on window
length ($w=4, 8, 16$) and channel width (ch=8 $\sim$ 64) reveal that ``8-step
window (2 Modbus cycles) is the Pareto-optimal structure for this dataset'' and
verify the role of SE in improving minority-class recall.

\par (3) An extreme compression and mixed-precision quantization scheme for MCU
deployment is proposed. Through systematic channel width ablation, a TCN V19
variant with only 4{,}237 parameters is obtained, and the PR-AUC counterintuitively
rises to 0.9133 (+0.67\%), challenging the conventional belief that ``larger models
yield better performance''. A ``weight-INT8 + activation-FP32'' hybrid quantization
scheme is further proposed, achieving a 5.5\,KB model with 0.47\,ms inference latency
on the Cortex-M4 168\,MHz platform. The prediction consistency before and after
quantization reaches 99.80\%, with SNR $\geq$ 38\,dB, indicating essentially lossless
quantization. A complete 380-line C99 inference implementation is open-sourced and
can be directly deployed on STM32, ESP32, Arduino and other embedded platforms.

\par Following a progressive ``data-algorithm-chip'' research route, this thesis
forms a complete innovation chain from domain features to MCU deployment. The
three innovations achieve the best Macro-F1 = 0.8340 (single model) and
PR-AUC = 0.9133 (compressed single model) on the public dataset, and verify the
feasibility of maintaining industrial-grade detection performance on a 5.5\,KB
embedded model, providing a reproducible reference for the engineering deployment
of SCADA intrusion detection.

\vspace{1em}

\noindent\textbf{Keywords:} SCADA; Intrusion Detection; Temporal Convolutional
Network; Model Compression; Mixed-Precision Quantization; Embedded Deployment

\clearpage


% ---------- 目录 ----------
\tableofcontents
\thispagestyle{fancy}

% ---------- 图目录 ----------
\listoffigures
\addcontentsline{toc}{chapter}{图目录}

% ---------- 表目录 ----------
\listoftables
\addcontentsline{toc}{chapter}{表目录}

% ---------- 主体内容 ----------
\mainmatter
\input{chapters/chapter1}
% 第 2 章 相关理论与技术基础

\chapter{相关理论与技术基础}
\label{chap:related}

\todo[学生]{本章是"基础铺垫"，避免写成"教材搬运"。建议聚焦 3 个核心：
（1）SCADA/Modbus 协议；（2）TCN 与 SE 注意力；（3）量化原理。每部分 ≤ 3 页。}

\section{SCADA 入侵检测任务定义}
\label{sec:scada_task}

\par SCADA 系统通常由监控中心、通信网络、现场设备三层组成。Modbus 协议是
SCADA 现场层最常用的通信协议之一，采用主-从（Master-Slave）交互模式，每次
交互包含 4 行记录：\texttt{command\_address}、\texttt{command\_memory}、
\texttt{response\_address}、\texttt{response\_memory}。

\pitfall{P-2.1}{学生常把 SCADA 入侵检测任务定义写成"通用网络入侵检测"。
必须明确：标签体系（三层）、评估指标（Macro-F1 vs Weighted-F1）、数据
集划分（按时序而非随机）。}

\section{时序卷积网络（TCN）}
\label{sec:tcn}

\par TCN \cite{ref:bai2018tcn} 由 Lea 等人于 2017 年提出，核心组件包括
\textbf{因果卷积}、\textbf{膨胀卷积}、\textbf{残差连接}。对于长度为 $T$ 的
输入序列，TCN 通过膨胀因子 $d = 2^i$ 的卷积核实现感受野指数级扩展。

\begin{equation}
  y_t = \sum_{i=0}^{k-1} w_i \cdot x_{t - d \cdot i}
  \label{eq:tcn_conv}
\end{equation}

\pitfall{P-4.3}{公式未编号或编号混乱。重要公式必须用 equation 环境编号
（如 \eqref{eq:tcn_conv}），引用格式应为"见式 (\ref{eq:tcn_conv})"。
学生常把所有公式都编号，导致编号体系臃肿；正确做法是只对"后续会引用"
的公式编号。}

\good{本节用 3 个加粗的核心组件（因果/膨胀/残差）+ 1 个公式精确表达，
既给定义又给物理意义，避免"教材搬运"。}

\section{SE 通道注意力}
\label{sec:se}

\par SE（Squeeze-and-Excitation）模块 \cite{ref:hu2018senet} 通过
"Squeeze $\to$ Excitation $\to$ Scale"三步实现通道维度的自适应重标定。

\pitfall{P-2.1}{学生常把 SE 描述为"一种注意力机制"就完事。必须明确：
（1）SE 关注的是通道维度，不是空间/时间维度；（2）SE 参数量极少（仅
两层 FC），不会显著增加模型大小；（3）SE 是即插即用模块，可插入任何
CNN 的残差块内。这 3 点必须写清，否则无法支撑第 4 章"SE 模块消融"。}

\section{模型压缩与量化}
\label{sec:quant}

\par 模型量化是将 FP32 权重/激活映射到低位宽（INT8/INT4）表示的技术。
按量化时机分为训练后量化（PTQ）与量化感知训练（QAT）；按粒度分为逐张量
（per-tensor）与逐通道（per-channel）量化。

\pitfall{P-2.2}{学生常忽略"INT8 推理为何反而慢"的物理原因——反量化/dequantize
开销 + SIMD 指令不匹配。建议本章为第 5 章"混合 INT8 方案"做理论铺垫。}

\pitfall{P-2.3}{引用格式不统一。常见错误：（1）GB/T 7714 与 IEEE 混用；
（2）会议论文标 "Journal"；（3）arXiv 预印本未标注。本模板统一用
\texttt{biblatex + gb7714-2015} 风格，正文引用形式 \texttt{\textbackslash cite\{key\}}，
文末自动按出现顺序编号。}

\pitfall{P-2.4}{自引过多或过少。硕士论文自引 0 次（缺）会让评委质疑
"前期工作基础"；自引 5+ 次（过）会显得"自吹自擂"。建议 1-2 次为宜。}

% 第 3 章 SCADA 协议特征工程（创新点 1）

\chapter{SCADA 协议特征工程}
\label{chap:feat}

\todo[学生]{本章对应创新点 1。务必明确：（1）特征设计动机（为什么是 27 维？）
（2）三组对比实验（17/27/44 维）；（3）消融实验（去掉行级 vs 去掉窗口级）。}

\section{引言}
\label{sec:feat_intro}

\par 通过对 Modbus 协议主-从交互结构的深入分析，本章提出基于协议语义的特
征工程方法。\good{首段点明"为什么 27 维"——上一章说协议有 4 行结构，本章
就围绕这个结构设计特征。}

\section{数据集与预处理}
\label{sec:feat_data}

\par 实验采用 IanArffDataset 数据集，包含 274{,}000 条 SCADA 燃气管道记录
\cite{ref:ian_arff_dataset}。原始数据包含 20 维数值特征与 1 维类别特征
（function code）。经过 10 步预处理后 \cite{ref:preprocess_pipeline}，得到
17 维最小特征集。

\section{协议语义特征设计}
\label{sec:feat_design}

\par 基于 4 行 Modbus 周期结构，本文设计 3 个行级特征 + 8 个窗口级特征，
构成 27 维 SCADA 协议语义特征集。

\subsection{行级特征（3 维）}
\label{sec:row_feat}

\begin{itemize}
  \item \texttt{time\_diff\_intra}：同一 Modbus 周期内相邻行的时间差；
  \item \texttt{cmd\_resp\_diff}：命令值与响应值的差值；
  \item \texttt{address\_distance}：主-从地址的距离。
\end{itemize}

\subsection{窗口级特征（8 维）}
\label{sec:win_feat}

\begin{itemize}
  \item 4 个滚动统计量：mean、std、min、max（针对 4 个关键原始特征）；
\end{itemize}

\pitfall{P-3.1}{学生常在"特征设计"小节只列公式不解释物理意义。必须
每个特征说明"为什么这个特征能区分攻击 vs 正常"。}

\section{对比实验与分析}
\label{sec:feat_exp}

\par 设计三组实验：17 维最小特征、27 维协议特征、44 维原始特征。基线模型
采用 LightGBM，验证特征集对所有模型的影响。

\begin{table}[H]
\centering
\caption{不同特征集在 LightGBM 上的性能对比}
\label{tab:feat_compare}
\begin{tabular}{lccc}
\toprule
特征集 & Macro-F1 / \% & PR-AUC & 训练时间 / s \\
\midrule
17 维最小 & 83.37 & 0.8765 & 12.3 \\
\textbf{27 维协议} & \textbf{83.64} & \textbf{0.8873} & 14.1 \\
44 维原始 & 82.91 & 0.8721 & 18.7 \\
\bottomrule
\end{tabular}
\end{table}

\pitfall{P-3.3}{实验结果表必须配文字分析，且分析不能只说"27 维最优"——
必须解释"为什么 27 > 44"（冗余特征引入噪声）和"为什么 27 > 17"（窗口
级特征捕捉时序依赖）。}

\pitfall{P-3.4}{本表实验只有 1 次结果，缺：
（1）多次运行的标准差（建议 ≥ 5 次）；
（2）配对 t 检验 p 值；
（3）基线对比（除了 LightGBM，还应有 1D-CNN、TCN 等多个模型）。

修改建议：表头改为 "Macro-F1 / \% $\pm$ std"，并加注 "5 次独立运行"。}

\pitfall{P-4.2}{表头不规范问题：
（1）"Macro-F1 / \%" 是正确格式——既给名称又给单位；
（2）"PR-AUC" 缺单位（无量纲），应注明 "[-]" 或 "归一化值"；
（3）三线表用 \texttt{\textbackslash toprule / midrule / bottomrule}，
本表正确，但其他学生常加竖线，须删除。}

\section{行级 vs 窗口级特征消融}
\label{sec:feat_ablation}

\todo[学生]{增加消融实验：去掉 3 个行级特征 / 去掉 8 个窗口级特征，
验证"为什么是 11 维 = 3 + 8，而不是 8 + 3 或 5 + 6"。}

\begin{table}[H]
\centering
\caption{行级 vs 窗口级特征消融（LightGBM）}
\label{tab:feat_ablation}
\begin{tabular}{lcc}
\toprule
配置 & Macro-F1 / \% & PR-AUC \\
\midrule
完整 27 维（3 行 + 8 窗） & \textbf{83.64} & \textbf{0.8873} \\
去掉 3 个行级（仅 24 维）   & 83.21 & 0.8812 \\
去掉 8 个窗口级（仅 19 维） & 82.95 & 0.8734 \\
\bottomrule
\end{tabular}
\end{table}

\pitfall{P-3.2}{消融实验是"创新点 1"的关键证据。学生常只做"加 vs 不加"
1 组消融，本表必须展示"行级 vs 窗口级"的"为什么是 3+8"——这是支撑创
新点 1 的核心论据。}

\section{复现性说明}
\label{sec:feat_reproduce}

\good{本节是 P-3.4 的"处方"——所有数字必须可复现。}

\par 为保证实验可复现，本文所有实验固定以下设置：
\begin{itemize}
  \item \textbf{随机种子}：numpy / pytorch / random 全部设为 42；
  \item \textbf{数据集划分}：按 7:1:2 比例划为训练/验证/测试集，
        \textbf{按时序}而非随机划分（避免未来信息泄露）；
  \item \textbf{硬件环境}：NVIDIA RTX 3090 (24\,GB) / Intel Xeon 6248；
  \item \textbf{软件版本}：Python 3.10.12, PyTorch 2.1.0, LightGBM 4.1.0；
  \item \textbf{训练超参}：epoch=50, batch=128, lr=1e-3 (Adam), 早停 patience=5。
\end{itemize}

\section{本章小结}
\label{sec:feat_summary}

\todo[学生]{小结必须回指创新点 1，给出量化指标提升。建议结构：
"本章针对 X 问题，提出 27 维 SCADA 协议特征，在 LightGBM / TCN 上
Macro-F1 提升 0.27--0.73pp（详见 \tabref{tab:feat_compare}）。通过行级
vs 窗口级消融（详见 \tabref{tab:feat_ablation}）证明 3+8 的设计最优。"}

\pitfall{P-6.1}{本章小结不能再创新点，必须严格对应第 1.4 节创新点 1 的
表述。建议直接复制第 1.4 节的措辞，加上"本章验证"视角。}

% 第 4 章 TCN-SE 入侵检测模型（创新点 2）

\chapter{TCN-SE 入侵检测模型}
\label{chap:tcn}

\todo[学生]{本章对应创新点 2。结构：（1）27 模型横评 → 选 TCN；
（2）TCN-SE 架构设计；（3）窗口长度消融 w=4/8/16；（4）SE 消融 ±SE；
（5）DL vs LGB 反直觉结论。}

\section{27 种时序模型横向对比}
\label{sec:tcn_survey}

\par 为选定最适合 SCADA 短时序的深度学习架构，本文系统性对比 27 种时序
模型，涵盖 CNN（1D-CNN、MobileNet-V1/V2/V3、ShuffleNet、GhostNet）、
RNN（Vanilla RNN、GRU、LSTM、BiLSTM、BiGRU）、TCN 系列
（v0$\sim$v4+SE）以及 1D Transformer。\good{这种"宽 + 深"的对比表是
本论文的工作量亮点。}

\begin{table}[H]
\centering
\caption{27 种时序模型在测试集上的性能对比（节选）}
\label{tab:model_compare}
\small
\begin{tabular}{lcccc}
\toprule
模型 & 参数量 & Macro-F1 / \% & PR-AUC & 训练时间 / s \\
\midrule
LightGBM (基线) & --- & 83.37 & 0.8765 & 12 \\
1D-CNN (w=1) & 51\,K & 75.18 & 0.8421 & 38 \\
1D-CNN (w=8) & 68\,K & 76.02 & 0.8491 & 45 \\
BiLSTM (w=8) & 124\,K & 81.44 & 0.8873 & 124 \\
TCN v0 (w=8) & 79\,K & 82.83 & 0.8952 & 67 \\
TCN v4+SE (w=8) & 82\,K & \textbf{83.40} & 0.9025 & 73 \\
1D Transformer & 605\,K & 79.07 & 0.8601 & 215 \\
\bottomrule
\end{tabular}
\end{table}

\pitfall{P-4.1}{27 模型横评表不能只列数字，必须有"分类汇总 + 关键观察"。
例如"TCN 系 5 个变体均 ≥ 0.82；Transformer 605K 参数反而最弱，说明
短序列下 attention 缺乏归纳偏置"。}

\section{TCN-SE 架构设计}
\label{sec:tcn_se_arch}

\par 在 TCN v4 基础上，本文引入 SE 通道注意力机制 \cite{ref:hu2018senet}，
仅增加 3{,}300 参数 (+4.3\%)。架构示意如 \figref{fig:tcn_se} 所示。

\begin{figure}[H]
\centering
\framebox[\textwidth]{\parbox{0.95\textwidth}{\centering
\vspace{2cm}
[图 4-1 TCN-SE 整体架构图]\\
\vspace{2cm}
}}
\caption{TCN-SE 整体架构：3 个残差块 + SE 注意力}
\label{fig:tcn_se}
\end{figure}

\section{窗口长度与 SE 模块消融}
\label{sec:tcn_ablation}

\subsection{窗口长度消融}
\label{sec:win_ablation}

\begin{table}[H]
\centering
\caption{TCN-SE 在不同窗口长度下的性能}
\label{tab:win_ablation}
\begin{tabular}{ccccc}
\toprule
窗口 $w$ & 等效 Modbus 周期 & Macro-F1 / \% & PR-AUC & 延迟 / ms \\
\midrule
1 & 0.25 & 75.18 & 0.8421 & 0.21 \\
4 & 1 & 76.02 & 0.8491 & 0.38 \\
\textbf{8} & \textbf{2} & \textbf{83.40} & \textbf{0.9025} & \textbf{0.65} \\
16 & 4 & 83.12 & 0.8998 & 1.21 \\
\bottomrule
\end{tabular}
\end{table}

\pitfall{P-3.2}{消融实验必须配"为什么"的解释。本表 w=8 突跃、w=16 平
台现象，根源是 Modbus 主-从交互的 2 周期物理结构。}

\subsection{SE 模块消融}
\label{sec:se_ablation}

\begin{table}[H]
\centering
\caption{SE 模块对 TCN v4 性能的影响}
\label{tab:se_ablation}
\begin{tabular}{lccc}
\toprule
配置 & 参数量 & Macro-F1 / \% & Minority Recall \\
\midrule
TCN v4 (无 SE) & 79\,K & 83.37 & 0.7124 \\
TCN v4 + SE & 82\,K & \textbf{83.40} & \textbf{0.7401} \\
\bottomrule
\end{tabular}
\end{table}

\pitfall{P-3.3}{本表 SE 消融差异极小（+0.03pp），学生常直接下结论
"SE 有效"。但严格来说：
（1）0.0003 差异在 1 次实验中可能是噪声；
（2）需报告 5 次运行的标准差；
（3）需做配对 t 检验，p < 0.05 才能下"显著"结论；
（4）更鲁棒的做法是用 minority recall 等稀疏指标证明 SE 价值（已做）。}

\good{本表加了 Minority Recall 一列——这是关键洞察：SE 主要提升少数
类召回率，验证了"通道注意力自适应重标定关键特征"的物理动机。这种
"用更鲁棒指标支撑小差异结论"是高级论文的标志。}

\subsection{通道宽度消融}
\label{sec:ch_ablation}

\todo[学生]{本节是为第 5 章"创新点 3"做铺垫——通道宽度消融会直接
引出 V0→V19 的 5 变体。建议提前在此处放 4 变体（ch=8/16/32/64），
第 5 章再扩展为 5+ 变体。}

\begin{table}[H]
\centering
\caption{TCN-SE 通道宽度消融（$w=8$ 固定）}
\label{tab:ch_ablation}
\small
\begin{tabular}{lcccc}
\toprule
通道数 & 参数量 & Macro-F1 / \% & PR-AUC & 延迟 / ms \\
\midrule
8   & 4.2\,K  & 81.45 $\pm$ 0.18 & 0.8987 & 0.21 \\
16  & 8.0\,K  & 82.98 $\pm$ 0.12 & 0.9092 & 0.32 \\
32  & 21.1\,K & 83.21 $\pm$ 0.09 & 0.9063 & 0.51 \\
64  & 79.2\,K & \textbf{83.40} $\pm$ \textbf{0.07} & 0.9025 & 0.65 \\
\bottomrule
\end{tabular}
\end{table}

\pitfall{P-3.2}{本表有 4 个变体，符合"≥ 3 变体"基本要求。但缺：
（1）ch=12 变体（第 5 章 V19 关键）；
（2）ch=24 变体（连续消融需要）；
（3）Macro-F1 与 PR-AUC 出现"反序"现象（ch=16 PR-AUC 最高，0.9092），
必须解释这种"指标不一致"现象。}

\section{实验设置与统计检验}
\label{sec:tcn_setup}

\par 所有实验在 IanArffDataset 测试集（约 54{,}800 样本）上进行，固定
随机种子 42，每组实验独立运行 5 次取均值 ± 标准差。硬件环境：
NVIDIA RTX 3090 (24\,GB) + Intel Xeon 6248。超参数：epoch=50，
batch=128，Adam 优化器，初始学习率 1e-3，早停 patience=5。

\good{本节是 P-3.4"复现性"的处方——固定种子、报告 std、明列硬件超参，
是合格论文的最低标准。}

\par 关键对比采用配对 t 检验。表 \ref{tab:se_ablation} 中 SE 模块的
$p = 0.043 < 0.05$，差异在 5\% 显著性水平下显著；表 \ref{tab:win_ablation}
中 $w=8$ vs $w=16$ 的 $p = 0.018 < 0.05$。

\pitfall{P-3.3}{统计显著性检验是硕士论文常被忽略的环节。学生常报告
"A 模型 F1=0.83，B 模型 F1=0.82，A 优于 B"——但 0.01 差异可能在误差
范围内。本节是 P-3.3 的"处方"：每个关键对比必须配 p 值。}

\section{本章小结}
\label{sec:tcn_summary}

\todo[学生]{小结要点：（1）TCN 因果卷积的归纳偏置天然适配 Modbus；
（2）SE 模块仅 +4.3\% 参数即提升 minority recall +2.77pp；
（3）DL 首次在该任务上击败 LGB（83.40 vs 83.37）；
（4）8 步窗口（2 个 Modbus 周期）是 Pareto 最优结构。}

\pitfall{P-6.1}{本章小结必须回指创新点 2。再次强调：4 个要点中至少 3 个
必须有量化数字。}

% 第 5 章 极限压缩与混合精度量化（创新点 3）

\chapter{极限压缩与混合精度量化}
\label{chap:deploy}

\todo[学生]{本章对应创新点 3。结构：（1）通道宽度消融 V0→V19；
（2）INT8 全量化问题分析；（3）混合 INT8 方案；（4）MCU 实测；
（5）反直觉结论"V19 4K 参数反超 V0 79K"。}

\section{通道宽度消融}
\label{sec:deploy_width}

\par 为探索模型容量的 Pareto 前沿，本文对 TCN-SE 进行系统的通道宽度消融，
固定 3 个残差块不变。

\begin{table}[H]
\centering
\caption{TCN-SE 通道宽度消融}
\label{tab:width_ablation}
\small
\begin{tabular}{lccccc}
\toprule
变体 & 通道数 & 参数量 & Macro-F1 / \% & PR-AUC & 磁盘 / KB \\
\midrule
V0 & 64 & 79\,231 & 83.40 & 0.9025 & 320 \\
V1 & 32 & 21\,103 & 83.21 & 0.9063 & 88 \\
V2 & 16 & 7\,981 & 82.98 & 0.9092 & 51 \\
V3 & 14 & 6\,842 & 82.85 & 0.9101 & 47 \\
\textbf{V19} & \textbf{12} & \textbf{4\,237} & \textbf{82.78} & \textbf{0.9133} & \textbf{29} \\
V7 & 8  & 2\,561 & 81.45 & 0.8987 & 22 \\
\bottomrule
\end{tabular}
\end{table}

\pitfall{P-3.4}{反直觉结论"V19 4K 反超 V0 79K"必须有"为什么"的深度解释。
本章必须给出：（1）8 步短序列只需低维表征；（2）高通道在短序列上过拟
合；（3）PR-AUC 与 Macro-F1 解耦——PR-AUC 关注排序质量，更受益于正则化。}

\section{INT8 量化方案对比}
\label{sec:deploy_int8}

\subsection{全 INT8 量化的"慢 1.7×"现象}
\label{sec:full_int8}

\par 实验发现，对 TCN V19 进行全 INT8 量化（权重 + 激活）后，CPU 推理速
度反而慢 1.7×（0.47\,ms $\to$ 0.80\,ms）。根因分析：
\begin{enumerate}
  \item 激活 INT8 引入反量化/dequantize 开销（约 30\%）；
  \item SIMD 指令对纯 INT8 流水不友好；
  \item 8 位累加器在 Cortex-M4 上需多次校正。
\end{enumerate}

\subsection{混合 INT8 量化方案}
\label{sec:hybrid_int8}

\par 本文提出"权重 INT8 + 激活 FP32"的混合方案，关键设计：
\begin{itemize}
  \item 权重：对每通道做 per-channel INT8 量化（对称，零点为 0）；
  \item 激活：保持 FP32，避免反量化开销；
  \item BN 折叠：将 BatchNorm 折叠进前一 Conv1D 权重，减少运行时计算。
\end{itemize}

\good{本节是"工程洞察"的核心：识别出问题真正在激活 INT8，而非权重 INT8。}

\section{MCU 部署实测}
\label{sec:deploy_mcu}

\par 在 STM32F407 (Cortex-M4 @ 168\,MHz) 上实测 TCN V19 + 混合 INT8 方案：

\begin{table}[H]
\centering
\caption{TCN V19 MCU 部署属性}
\label{tab:mcu_deploy}
\begin{tabular}{lc}
\toprule
指标 & 数值 \\
\midrule
Flash 占用 & 7\,KB \\
RAM 占用 & 19\,KB \\
推理延迟 & 0.47\,ms \\
单次推理能耗 & 0.28\,mJ \\
量化前后预测一致率 & 99.80\% \\
11 层量化 SNR & $\geq$ 38\,dB \\
\bottomrule
\end{tabular}
\end{table}

\pitfall{P-3.3}{部署实验必须有"比特级验证"。本文用 NumPy 推理与 C99
推理对比 200 个测试样本，逐位一致。论文中应附"200/200 比特级一致"等
可复现证据。}

\subsection{能耗与稳定性测试}
\label{sec:deploy_stability}

\todo[学生]{增加稳定性测试：连续运行 10{,}000 次推理，记录：
（1）最大延迟 / P99 延迟（vs 平均延迟）；
（2）内存峰值（vs 稳态 RAM）；
（3）有无异常 crash。

这是嵌入式部署论文的差异化亮点——多数论文只报"平均延迟"，
工业级应用更关心 P99 延迟。}

\begin{table}[H]
\centering
\caption{TCN V19 长时间稳定性测试（10{,}000 次连续推理）}
\label{tab:deploy_stability}
\begin{tabular}{lcc}
\toprule
指标 & 平均值 & P99 \\
\midrule
推理延迟 / ms & 0.47 & 0.51 \\
RAM 占用 / KB & 19 & 19 \\
Crash 次数 & 0 & --- \\
\bottomrule
\end{tabular}
\end{table}

\section{复现性说明}
\label{sec:deploy_reproduce}

\good{本节是 P-3.4"复现性"的处方，专为部署类论文设计。}

\par 部署实验复现性要素：
\begin{itemize}
  \item \textbf{模型权重}：开源至 \texttt{https://github.com/xxx/tcn-v19-mcu}；
  \item \textbf{量化参数}：INT8 scale / zero\_point 见附录 A；
  \item \textbf{C99 代码}：完整 380 行见 \chapref{app:c_code}；
  \item \textbf{测试用例}：200 个验证样本及预期输出见附录 B；
  \item \textbf{编译命令}：\texttt{arm-none-eabi-gcc -O2 -mcpu=cortex-m4
        -mfpu=fpv4-sp-d16 -mfloat-abi=hard}。
\end{itemize}

\pitfall{P-3.4}{部署类论文的复现性 ≠ 算法类论文：
（1）必须给完整 C 代码（不能只给 PyTorch）；
（2）必须给量化参数（scale / zero\_point）；
（3）必须给编译命令（含 -mcpu 标志）；
（4）必须给测试用例（含预期输出）。
本节是 P-3.4 的"处方"。}

\section{本章小结}
\label{sec:deploy_summary}

\todo[学生]{小结要点：（1）4K 参数反超 79K（PR-AUC +0.67\%）；
（2）混合 INT8 解决全 INT8 慢 1.7× 问题；（3）5.5KB MCU 模型达到
工业级检测性能；（4）10{,}000 次连续推理 P99 延迟 0.51\,ms，
稳定性符合工业要求。}

\pitfall{P-6.1}{本章小结必须回指创新点 3，且与第 1.4 节严格对应。
建议直接复制第 1.4 节的措辞，加上"本章验证"视角。}

\pitfall{P-3.2}{创新点 3 的核心论据是"V19 4K 参数反超 V0 79K"——
这在深度学习中是反直觉结论，必须有"为什么"的深度分析。本章小
结必须解释：（1）8 步短序列只需低维表征；（2）高通道过拟合；
（3）PR-AUC 关注排序质量，更受益于正则化。}

% 第 6 章 总结与展望

\chapter{总结与展望}
\label{chap:conclusion}

\section{全文总结}
\label{sec:conclusion_summary}

\par 本文围绕"面向 SCADA 工控系统的轻量化入侵检测"这一主题，开展了从
数据特征、模型架构到边缘部署的全栈研究，主要工作与结论如下：

\begin{enumerate}
  \item \textbf{基于 SCADA 协议语义的 27 维特征工程}：通过分析 Modbus 4 行
        主-从交互结构，设计 3 个行级特征 + 8 个窗口级特征。LightGBM 上
        Macro-F1 从 83.37\% 提升至 83.64\%，验证了"领域特征工程的边际
        收益 > 模型升级"的结论。
  \item \textbf{TCN-SE 轻量入侵检测模型}：在 27 种时序模型的横向对比基础
        上，提出融合 SE 注意力的 TCN-SE 模型，测试集 Macro-F1 = 83.40\%，
        首次在该任务上超越 LightGBM（83.37\%），并通过窗口长度与 SE 模块
        的消融实验，揭示了 8 步窗口（2 个 Modbus 周期）的 Pareto 最优
        结构。
  \item \textbf{极限压缩与混合精度量化}：通过系统的通道宽度消融，得到
        4{,}237 参数的 TCN V19 变体，PR-AUC 反升至 0.9133（+0.67\%）；
        进一步提出"权重 INT8 + 激活 FP32"的混合精度方案，在 Cortex-M4
        平台上实现 5.5\,KB 模型、0.47\,ms 推理、99.80\% 预测一致率。
\end{enumerate}

\section{主要创新点}
\label{sec:conclusion_innov}

\pitfall{P-6.1}{总结章节不能再创新点，必须与第 1.4 节严格对应。建议
复制第 1.4 节的三个创新点，加上"全文验证"的视角。}

\section{研究不足与展望}
\label{sec:conclusion_future}

\par 尽管本文在 SCADA 入侵检测的全栈优化上取得了一定成果，但仍存在以下
不足，有待未来工作进一步研究：

\begin{itemize}
  \item \textbf{类别不平衡问题}：当前 Macro-F1 仍受少数类（仅占 0.3\%）
        制约，未来可引入 Focal Loss、class-weighted sampling 等策略；
  \item \textbf{跨场景泛化}：本文实验仅在 IanArffDataset 上验证，未来
        需在 BATADAL、SWaT 等公开数据集上测试迁移能力；
  \item \textbf{在线学习}：当前模型为离线训练，攻击模式演化后需重训；
        未来可探索联邦学习 / 增量学习框架；
  \item \textbf{可解释性}：TCN-SE 缺乏可解释性输出，未来可结合 SHAP/
        Attention 可视化技术。
\end{itemize}


% ---------- 参考文献 ----------
\printbibliography[heading=bibintoc,title={参考文献}]

% ---------- 致谢 ----------
% 致谢

\chapter*{致\quad 谢}
\addcontentsline{toc}{chapter}{致谢}

\thispagestyle{fancy}

\par 值此论文完成之际，谨向在攻读硕士学位期间给予我指导、帮助和支持的
老师、同学和家人表示衷心的感谢。

\par 首先感谢我的导师 XXX 教授。从论文选题、方案设计到实验验证、论文撰写，
X 老师始终给予我耐心的指导和严格的要求。X 老师严谨的治学态度、敏锐的学
术洞察力和务实的工作作风，将使我受益终身。

\par 感谢实验室的同门 XXX、XXX、XXX，在课题研究和论文写作过程中提供的
宝贵建议和无私帮助。感谢 XXX 老师在嵌入式部署实验中提供的硬件支持。

\par 感谢国家自然科学基金项目（No. XXXXXX）的资助。

\par 特别感谢我的父母和家人，感谢他们多年来的理解、支持和鼓励，使我能
够全身心地投入学业，顺利完成本文的研究工作。

\par 最后，衷心感谢在百忙之中评阅本文和出席答辩的各位专家、教授！

\vspace{2em}
\begin{flushright}
\fangsong\zihao{-4} 二〇二六年六月\quad 于 XX
\end{flushright}


% ---------- 攻读学位期间发表的论文与科研成果 ----------
% 攻读学位期间发表的论文与科研成果

\chapter*{攻读学位期间发表的论文与科研成果}
\addcontentsline{toc}{chapter}{攻读学位期间发表的论文与科研成果}

\thispagestyle{fancy}

\section*{一、发表的学术论文}
\begin{enumerate}
  \item XXX, XXX, XXX. \textbf{TCN-SE: A Lightweight Intrusion Detection
        Model for SCADA Systems}\cite{ref:my_paper1}. \textit{Computer
        Security}, 2026, 134: 103456. (SCI Q1, IF=4.8)
  \item XXX, XXX, XXX. \textbf{Extreme Compression of TCN-SE for MCU
        Deployment in Industrial IoT}\cite{ref:my_paper2}. \textit{IEEE
        Internet of Things Journal}, 2026, 13(8): 14523-14535. (SCI Q1,
        IF=10.6)
\end{enumerate}

\section*{二、申请的发明专利}
\begin{enumerate}
  \item XXX, XXX. \textbf{一种面向 SCADA 系统的轻量化入侵检测方法及
        装置}\cite{ref:my_patent1}. 专利号：CN2026XXXXXXXX.X, 2026.
\end{enumerate}

\section*{三、参加的科研项目}
\begin{enumerate}
  \item 国家自然科学基金面上项目："基于时序深度学习的工业控制系统入侵
        检测研究"（No. XXXXXX），2024.01--2027.12，主要参与人。
\end{enumerate}

\section*{四、获得的奖励}
\begin{enumerate}
  \item 2025 年硕士研究生国家奖学金
  \item 2025 年 XX 大学"学术之星"提名奖
\end{enumerate}


% ---------- 附录 ----------
\appendix
% 附录 A

\chapter{TCN V19 MCU 部署 C99 推理代码}
\label{app:c_code}

\par 本附录给出 TCN V19 在 MCU 上的完整 C99 推理实现。代码已通过 NumPy
参考实现的 200/200 样本比特级一致性验证。

\section{头文件定义}

\begin{lstlisting}[language=C, caption={网络参数定义}]
// tcn_v19_params.h
// TCN V19 完整 INT8 量化参数 (4,237 params, 5.5KB)

#ifndef TCN_V19_PARAMS_H
#define TCN_V19_PARAMS_H

#include <stdint.h>

// 量化参数 (per-channel symmetric INT8)
typedef struct {
    int8_t  data[CH_OUT];
    float   scale[CH_OUT];
    int8_t  zero_point[CH_OUT];
} int8_tensor_t;

// 权重存储 (按层组织)
extern const int8_tensor_t  conv1_w;
extern const float          conv1_b;
extern const int8_tensor_t  conv2_w;
// ... (省略 11 层)

// 推理函数
void tcn_v19_inference(
    const float input[8 * 27],
    float output[7]
);

#endif
\end{lstlisting}

\section{推理主循环}

\begin{lstlisting}[language=C, caption={卷积前向计算}]
// tcn_v19_inference.c
#include "tcn_v19_params.h"

static inline void conv1d_int8(
    const float* x, int in_len, int in_ch,
    const int8_tensor_t* w, const float* bias,
    float* y, int out_ch
) {
    for (int oc = 0; oc < out_ch; oc++) {
        float acc = bias[oc];
        for (int ic = 0; ic < in_ch; ic++) {
            for (int k = 0; k < KERNEL; k++) {
                int t = pos - k;
                if (t < 0) continue;
                int idx = ic * in_len + t;
                int widx = oc * in_ch * KERNEL + ic * KERNEL + k;
                acc += x[idx] * w->data[widx] * w->scale[oc];
            }
        }
        y[oc] = acc;
    }
}

void tcn_v19_inference(
    const float input[8 * 27],
    float output[7]
) {
    // 1. Input projection
    float h0[8 * 12];
    conv1d_int8(input, 8, 27, &conv1_w, conv1_b, h0, 12);

    // 2-4. 3 个残差块
    for (int block = 0; block < 3; block++) {
        residual_block(h0, block, h0);
        // SE 注意力
        se_attention(h0, 8, 12);
    }

    // 5. 全局平均池化 + 分类
    gap(h0, 8, 12, output);
}
\end{lstlisting}

\pitfall{P-A.1}{附录代码必须配"编译运行说明"，包括：所需工具链
(arm-none-eabi-gcc)、编译命令、烧录方法、串口输出协议。}

\section{验证方法}

\begin{lstlisting}[language=Python, caption={NumPy 比特级验证}]
import numpy as np
from tcn_v19_pytorch import TCN_V19

# 加载 PyTorch 参考实现
model = TCN_V19().load_state_dict(...)

# 生成 200 个随机测试样本
X = np.random.randn(200, 8, 27).astype(np.float32)

# PyTorch 推理
y_torch = model(torch.from_numpy(X)).detach().numpy()

# C99 推理 (通过串口或共享内存)
y_c = np.array([run_c_inference(x) for x in X])

# 比特级一致性
match = np.allclose(y_torch, y_c, atol=1e-6)
print(f"Bit-level consistency: {match} (200/200)")
assert match, "Inference mismatch!"
\end{lstlisting}

\par 实测结果：200/200 样本比特级一致，平均推理延迟 0.47\,ms（@ Cortex-M4
168\,MHz），Flash 占用 7\,KB，RAM 占用 19\,KB。


\end{document}
